\documentclass{article}
\usepackage{ctex}

\begin{document} % begin{document} 后面不加代码默认空两格 \noindent 可以消除空格
第一回\hspace{1em}灵根育孕源流出　心性修持大道生 \\ % 换行
% 1em空一个字的宽度 ，还可以1cm，1in（英寸）
% 用\hspace时前后空格也影响宽度

% 诗曰： \newline混沌未分天地乱 错误。\newline无大括号，后面要空一格
诗曰： \newline % 换行
\vspace{1em} % 垂直空格
混沌未分天地乱，茫茫渺渺无人见。 \\
自从盘古破鸿蒙，开辟从兹清浊辨。 \\
覆载群生仰至仁，发明万物皆成善。 \\
欲知造化会元功，须看西游释厄传。 \\
% 空段之后会缩进两个字节
  盖闻天地之数，有十二万九千六百岁为一元。将一元分为十二会，乃子、丑、寅、卯、辰、巳、午、未、申、酉、戌、亥之十二支也。每会该一万八百岁。且就一日而论：子时得阳气，而丑则鸡鸣；寅不通光，而卯则日出；辰时食后，而巳则挨排；日午天中，而未则西蹉；申时晡而日落酉；戌黄昏而人定亥。譬于大数，若到戌会之终，则天地昏蒙而万物否矣。再去五千四百岁，交亥会之初，则当黑暗，而两间人物俱无矣，故曰混沌。又五千四百岁，亥会将终，贞下起元，近子之会，而复逐渐开明。邵康节曰：“冬至子之半，天心无改移。一阳初动处，万物未生时。”到此，天始有根。 \par % 同样会分段并缩进2格

\end{document}